\clearpage\thispagestyle{empty}
\pdfbookmark[0]{Quellen zu Pallas (31. August 1813)}{pallas1813aug-sources}
\begin{center}\textbf{Quellen zu Pallas (31. August 1813)}\end{center}
\noindent Diese KI-Übertragung ins heutige Deutsch folgt dem Brief vom 31. August 1813 in den \href{https://gdz.sub.uni-goettingen.de/id/PPN236007467}{\textit{Werken}, Band~VI} von Gauß (Göttingen, 1874), S.~368--369. Er erschien zuerst als Nr.~XXI, \textit{Aus einem Schreiben des Hrn. Prof. Gauss}, in der Zeitschrift \textit{Monatliche Correspondenz zur Beförderung der Erd- und Himmelskunde}, Band~XXVIII (August 1813), S.~\href{https://api.digitale-sammlungen.de/iiif/image/v2/bsb10538617_00205/full/1600,/0/default.jpg}{197}, \href{https://api.digitale-sammlungen.de/iiif/image/v2/bsb10538617_00206/full/1600,/0/default.jpg}{198}. Die Links führen zum digitalisierten Exemplar der Bayerischen Staatsbibliothek.
\par\smallskip\noindent Der Erstdruck beginnt mit vier Punkten, die Werkausgabe mit drei Strichen. Für den 16. August stehen im Erstdruck in beiden Deklinationsfeldern fünf Punkte. Die Werkausgabe hat fünf in der Beobachtungstabelle und vier in der Vergleichstabelle. Sie lässt zudem die Wiederholungsstriche für den Monat in der ersten Tabelle weg und verwendet für die Zeiteinheiten hochgestellte u, m und s. KI-Übertragung und KI-Satz.
\clearpage
